% =====================================================================
\chapter{C3: evicting a compromised TUAV by quorum}\label{ch:c3}
% =====================================================================
\begin{keybox}[colframe=quorum,colback=quorum!5]
With the TA offline, nobody in the field can revoke a TUAV whose software has been taken over. In C3, \textbf{$t$ neighbouring TUAVs vote}: each
signs ``evict $T_{\mathrm{target}}$'', and $t$ distinct votes form an \textbf{eviction certificate} that every UAV and TUAV can check with public keys
alone. Members of the evicted TUAV stop trusting it and re-join a neighbour by the full join.
\end{keybox}

\section{Why it is needed}
After C1 the TUAV is the root of trust of its cell: it admits members and holds every pairwise key and the group key. A TUAV can be lost in two ways.
\textbf{Physical capture} -- the drone leaves the tether and falls into the wrong hands -- is handled by C4 (Chapter~\ref{ch:c4}). \textbf{Software
compromise while still tethered} -- a hijacked radio stack or a malicious update -- cannot be noticed by the tether, and the base paper has no revocation
mechanism for TUAVs at all. C3 covers this case without the TA.

\section{Votes and certificates}
TUAV $j$ votes with its ordinary C1 key (role label $\ep=0$):
\[
\mathit{vote}_j=\mathrm{Sign}_{\sk_j}\big(\texttt{evict-tuav}\cat\ID_{\mathrm{target}}\cat\ts_j\big)=(R_j,v_j).
\]
\begin{definition}[Valid eviction certificate]\label{def:cert}
A set of votes $\mathcal C$ is a valid certificate against $T_{\mathrm{target}}$ when all four hold:
\begin{enumerate}[leftmargin=*,label=\textbf{V\arabic*.}]
\item it contains votes from at least $t$ \emph{distinct} TUAV identities;
\item $T_{\mathrm{target}}$ is not among the signers;
\item every vote names the same target;
\item the weighted batch check~\eqref{eq:bv} holds over the votes, with $h_j=\Ho(\ID_j\cat0\cat X_j\cat Y_j)$.
\end{enumerate}
\end{definition}
\noindent Because TUAV credentials are bound to $\ep=0$ inside $h_j$, a member's credential can never cast a vote (role binding, fix A4).

\begin{figure}[H]
\centering
\resizebox{0.95\linewidth}{!}{%
\begin{tikzpicture}[font=\sffamily\small]
  \node[card=bad,minimum width=2.2cm] (t7) at (0,0) {$T_7$\\ \scriptsize compromised};
  \node[card=quorum] (t1) at (4.4,2.0) {$T_1$};
  \node[card=quorum] (t2) at (4.4,0) {$T_2$};
  \node[card=quorum] (t3) at (4.4,-2.0) {$T_3$};
  \draw[-{Stealth},thick,bad,dashed] (t7.east) -- node[lbl,above,sloped]{evidence} (t1.west);
  \draw[-{Stealth},thick,bad,dashed] (t7.east) -- (t2.west);
  \draw[-{Stealth},thick,bad,dashed] (t7.east) -- (t3.west);
  \foreach \n in {t1,t2,t3}{\draw[-{Stealth},thick,quorum] (\n.east) -- (8.1,0);}
  \node[card=quorum,text width=3.6cm] (cert) at (9.9,0) {\textbf{certificate}\\ $\ge t$ distinct votes\\ ``evict $T_7$''};
  \foreach \y in {1.0,0,-1.0}{\pic[color=uavc,scale=0.8] at (-3.4,\y) {drone};}
  \node[font=\scriptsize\sffamily\bfseries,text=uavc] at (-3.4,-1.8) {$T_7$'s members};
  \draw[-{Stealth},thick,quorum,dashed] (cert.south) |- (-3.4,-3.6) -- (-3.4,-2.15);
  \node[lbl,text=quorum] at (3.4,-3.95) {members check V1--V4, block $T_7$ for ever, re-join a neighbour by the full join (C1)};
  \draw[bad,line width=1.6pt] (-2.6,0.5) -- (-1.0,0.1);
  \draw[bad,line width=2pt] (-1.95,0.55) -- ++(0.3,-0.3) (-1.95,0.25) -- ++(0.3,0.3);
\end{tikzpicture}}
\caption{C3: three TUAVs see evidence, vote, and the certificate reaches the members of the compromised TUAV.}
\label{fig:c3vote}
\end{figure}

\section{The protocol}
\begin{stepbox}{quorum}{EV1--EV5 \;Eviction with the TA offline}
\steplabel{quorum}{EV1 (evidence).} A TUAV obtains evidence against a neighbour: \emph{cryptographic} evidence -- two conflicting messages signed by
the target (for instance two different rekeys with the same version $\nu$), or a signed message that breaks a rule (a decreasing version, a ticket for
an evicted member) -- or \emph{operational} evidence such as a crew or sensor report.\\
\steplabel{quorum}{EV2 (vote).} It signs $\mathit{vote}_j$ and sends it, with the evidence, to its neighbours.\\
\steplabel{quorum}{EV3 (assemble).} Any TUAV that holds $t$ distinct valid votes forms the certificate $\mathcal C$.\\
\steplabel{quorum}{EV4 (spread).} $\mathcal C$ is broadcast in the neighbouring cells and relayed into the target's cell.\\
\steplabel{uavc}{EV5 (act).} A member checks V1--V4, blocks the target \emph{for ever}, discards its keys, and re-joins a neighbouring TUAV by the full
join UA1--UA5 -- never with a C2 ticket from the evicted TUAV. Neighbouring TUAVs refuse its tickets and link offers.
\end{stepbox}

\begin{figure}[H]
\centering
\resizebox{0.92\linewidth}{!}{%
\begin{tikzpicture}[font=\sffamily\small,node distance=0.5cm]
  \node[fio=quorum,text width=2.6cm] (in) {receive a certificate $\mathcal C$};
  \node[fdec=quorum,right=of in] (v1) {V1: $\ge t$\\ distinct?};
  \node[fdec=quorum,right=of v1] (v2) {V2--V3:\\ target ok?};
  \node[fdec=quorum,right=of v2] (v4) {V4: batch\\ holds?};
  \node[fstart=good,right=of v4,text width=2.6cm] (ok) {block target;\\ full join to a neighbour};
  \node[fstart=bad,below=1.1cm of v2,text width=3.0cm] (no) {ignore $\mathcal C$\\ (keep the current TUAV)};
  \draw[flow] (in) -- (v1); \draw[flow] (v1) -- node[lbl,above]{yes} (v2); \draw[flow] (v2) -- node[lbl,above]{yes} (v4); \draw[flow] (v4) -- node[lbl,above]{yes} (ok);
  \draw[flow] (v1) |- node[lbl,pos=0.3,left]{no} (no); \draw[flow] (v2) -- node[lbl,right]{no} (no); \draw[flow] (v4) |- node[lbl,pos=0.3,right]{no} (no);
\end{tikzpicture}}
\caption{Flowchart: how a member checks an eviction certificate.}
\label{fig:c3flow}
\end{figure}

\section{Security}
\begin{theorem}[Eviction safety]\label{thm:evsafe}
Suppose at most $f\le t-1$ TUAVs are compromised, signatures are unforgeable (Theorem~\ref{thm:unforge}), and honest TUAVs vote only on evidence of
misbehaviour. Then, except with negligible probability, no valid certificate names an honest TUAV whenever the evidence is cryptographic; for operational
evidence the same holds as long as no honest voter is misled.
\end{theorem}
\begin{proof}
A valid certificate carries $t$ valid votes from distinct identities (V1, V4). At most $f<t$ of them can come from compromised TUAVs, and producing a valid
vote under an honest TUAV's key is a forgery. Hence at least one honest TUAV voted, which it does only with evidence against the target. Cryptographic
evidence against an honest TUAV would consist of messages carrying its valid signatures that it never produced -- again a forgery.
\end{proof}

\begin{theorem}[Eviction liveness]\label{thm:evlive}
If at least $t$ honest TUAVs obtain evidence against a TUAV and can reach its members, a valid certificate forms and every member that receives it stops
using that TUAV.
\end{theorem}
\begin{proof}
The $t$ honest votes name the same target and come from distinct identities other than the target, so V1--V3 hold; V4 holds by
Lemma~\ref{lem:c1corr}(ii). EV5 then applies.
\end{proof}

\begin{corollary}[Choosing $t$]
If up to $f$ of a TUAV's $n_T$ neighbours may be compromised, choose $t=f+1$ for safety; liveness then needs $n_T-f\ge t$, i.e.\ $n_T\ge2f+1$.
\end{corollary}

\begin{remark}[Relation to PBFT]
The pattern -- a quorum of signed statements that anyone can check -- is the view change of PBFT~\cite{castro99}, where a new leader must show $2f+1$
signed messages out of $3f+1$ replicas. PBFT needs the larger quorum because replicas must agree on an \emph{order} of requests. Eviction is a monotone
flag (``evicted'' never reverts) with no order to agree on, so a quorum of $f+1$ already prevents framing. C3 is a new \emph{use} of this pattern --
removing a key-management node with the authority offline -- not new cryptography.
\end{remark}

\section{Cost}
A vote is one signature. A member checks $t$ votes in one weighted batch, about $t+2$ scalar multiplications; for $t=2$ that is $4\times0.9018\approx3.6$\,ms
with the base paper's constants.

\begin{limitbox}
\begin{itemize}[leftmargin=*]
\item $t$ or more compromised TUAVs can still evict honest ones. C4 reduces this risk: a \emph{captured} drone cannot sign, so it cannot vote.
\item C3 decides \emph{whether} to evict from evidence; it does not detect a compromise by itself.
\item Members that cannot reach any honest neighbour cannot re-join until one comes into range.
\end{itemize}
\end{limitbox}
